\section{Related Work}
\label{sec:related}

\para{Benchmark contamination.} The risk of data contamination in LLM evaluation has been widely recognized. \citet{zhang2024gsm1k} created GSM1K, a human-authored mirror of \gsmk, and found that several open-source models (Phi, Mistral, Yi families) show 2--8\% performance gaps between the original and the mirror, indicating contamination. Crucially, frontier models showed no such gap, suggesting that sufficiently capable models are naturally contamination-resistant. Our work extends this by systematically measuring the contamination effect across multiple benchmarks rather than studying a single benchmark in isolation.

\para{Harder and expert-level benchmarks.} Several benchmarks have been designed to resist contamination through difficulty. \citet{rein2023gpqa} introduced \gpqa, where PhD non-experts with full internet access score only 34\%---barely above the 25\% random baseline. \citet{tiglab2024mmlupro} extended \mmlu with 10 answer choices and harder questions, reducing random guessing from 25\% to 10\%. \citet{phan2025hle} pushed this further with Humanity's Last Exam, where frontier models score below 14\%. Our experiments provide direct evidence that these design choices translate into measurable finetuning resistance.

\para{Dynamic and temporal benchmarks.} An orthogonal approach ensures freshness rather than difficulty. LiveBench~\citep{white2024livebench} replaces approximately one-sixth of its questions monthly using post-training-cutoff materials. DyVal~\citep{zhu2023dyval} generates evaluation instances procedurally from random directed acyclic graphs, achieving near-zero collision probability. \citet{zhu2024contamination_survey} survey over 20 dynamic benchmarks and propose six quality criteria (correctness, scalability, collision resistance, complexity stability, diversity, interpretability), finding that no single benchmark satisfies all six. Our work complements the dynamic approach by characterizing which \emph{static} benchmarks are most resistant, which is relevant because static benchmarks remain dominant in practice.

\para{Contamination-resistant design.} WinoGrande~\citep{sakaguchi2019winogrande} pioneered bias reduction through the AFLITE algorithm, which removes instances solvable by superficial patterns. \gsmsym~\citep{mirzadeh2025gsmsymbolic} randomizes numeric values in math problems to test whether models truly reason or memorize. ARC~\citep{chollet2019arc} uses visual abstract reasoning tasks designed to resist memorization, though recent work has shown that test-time training on synthetic data can still achieve high scores. Our finding that AFLITE-reduced \wino is actually the \emph{least} resistant benchmark challenges the assumption that statistical bias reduction prevents finetuning inflation.

\para{Measuring memorization.} Beyond benchmark design, several methods detect memorization directly. Per-character log-likelihood analysis~\citep{zhang2024gsm1k} correlates with overfitting on known benchmarks. Answer-only classifiers trained on \gpqa's answer choices cannot exceed chance~\citep{rein2023gpqa}, confirming no spurious features. Our approach differs by directly measuring the \emph{causal effect} of fine-tuning through controlled experiments rather than detecting memorization post-hoc.
